\subsubsection{混合机型的成本分担}
表\ref{tab:rev_q1_types}按机型分解18架次正式方案。B型与C型各9架次，但分别承担225 kg和533 kg。相同架次数并不表示相同任务质量或能源成本；混合机队按服务区需求和具体航段选择机型，而非平均分配物资。
\begin{table}[htbp]\centering\small
\caption{18架次正式组批方案的机型分工}\label{tab:rev_q1_types}
\begin{tabular}{lrrrr}\toprule
机型 & 架次 & 质量/kg & 能耗/kWh & 累计作业/min\\\midrule
A & 0 & 0 & 0.000000 & 0.000\\
B & 9 & 225 & 19.511736 & 255.512\\
C & 9 & 533 & 39.619560 & 290.755\\
\bottomrule\end{tabular}
\end{table}


仅C型精确方案同样达到18架次，因而可作为控制架次数的对照。混合方案减少14.803612 kWh能耗、17.712306 min累计作业时间，说明优化收益并非只是少飞几次，而来自同架次数下的机型与箱组选择。质量优先FFD也产生18架次，却没有获得相同的次级成本，进一步说明完整物理成本需要进入组合选择。

18与19架次方案之间的差分对应每节省1 kWh增加约284.86 min累计作业时间。这个比例用于解释两份离散方案的选择代价，反映S008由一项C型任务改成两项B型任务后，准备、交接和往返工作量的变化。基准架次优先规则选择18架次，19架次则保留为偏重能源目标的可执行备选。

\subsubsection{参数变化与实际任务调整}
对批型$p$定义$\rho_p=1-E_p/E_m$，其安全余量可行条件是$\alpha\le\rho_p$。因此候选集合仅在有限临界值处变化，连续参数能够通过不可拆组合产生离散决策跃迁。事件枚举据此覆盖整个参数范围，比仅比较若干网格点更能定位任务调整时机。

“最少架次数不变”和“原任务无需改变”应分别判断。例如28\%与32\%场景均为20架次，但机型构成不同，实际箱组仍需更新。基准原方案在23.08835\%以内保持可行，则意味着在这一安全要求范围内可以直接沿用已发出的任务清单。超过首个阈值时先复核S004，随后关注S008，使参数分析对应具体服务区与执行操作。
